\chapter{Introduction}
\label{chap:introduction}

Even though universities rely heavily on digital systems, most computer labs still run practical exams manually \cite{pmc_cbt}. Coordinators share schedules on WhatsApp. Students sign paper sheets for attendance and email their code files at the end. During the exam, a couple of invigilators have to watch dozens of screens at once. It is clear that this setup does not scale.

Practical exams are crucial for computer science and software engineering students. We cannot test coding skills, debugging, or system design with simple multiple-choice questions. Students need a real machine, an IDE, and a working compiler. At the same time, the invigilator needs a way to see what is happening on every computer without constantly pacing the room.

We designed SmartExam to bridge this gap. The platform combines scheduling, student identity checks, desktop lockdown, live telemetry, auto-saves, and automated grading into a single system \cite{pmc_cbt}.

The issue is not a lack of effort from invigilators. The real problem is that the tools they use do not talk to each other. Schedules are on one sheet, attendance is on another, and student files are scattered across different folders. When a student disputes their grade or someone is suspected of cheating, the teacher has no solid evidence to review.

SmartExam brings all these pieces together. Admins set up access roles from a central panel. Before the exam, teachers create questions, upload test cases, and whitelist permitted software. When students log in, the desktop client locks down their machine to the exam workspace. Invigilators watch a live dashboard grid showing window focus and active processes. Once the exam concludes, the system compiles grades and logs into a PDF report \cite{integ1}, creating a complete evidence trail.

SmartExam is that unified system.
